\documentclass[10pt,a4paper]{amsart}
\usepackage[margin=19mm]{geometry}
\usepackage{amsmath,amssymb,mathtools}
\usepackage[scr=rsfs,cal=euler]{mathalfa}
\usepackage{xfrac}
\usepackage[unicode,hidelinks,bookmarks=false]{hyperref}
\newcommand{\ie}{i.e.\ }
\newcommand{\cf}{cf.\ }
\newcommand{\resp}{respectively\ }
\newcommand{\Subsection}{\subsection}
\providecommand{\nobreakdash}{\nobreak\hbox{-}}
\setlength{\emergencystretch}{2em}
\multlinegap0pt
\usepackage{amssymb}
\usepackage{mathtools}
\usepackage{float}
\usepackage{tikz}
\usepackage{xcolor}
\usepackage{bm}
\usepackage[shortlabels]{enumitem}
\newtheorem{theorem}{Theorem}
\title{Pricing and Hedging of Discretely Monitored Asian Options\\in the Volterra--Heston Model}
\author{Gijs Custers, Sven Karbach, Martin Friesen}
\date{Source: arXiv:2609.07169v1 (CC BY 4.0)}
\begin{document}
\maketitle

\noindent\emph{Source and licence.} Pricing and Hedging of Discretely Monitored Asian Options\\in the Volterra--Heston Model. Version arXiv:2609.07169v1 (CC BY 4.0), distributed by arXiv under the Creative Commons Attribution 4.0 International licence, http://creativecommons.org/licenses/by/4.0/, as recorded in the arXiv licence metadata for this exact version and shown on the abstract page https://arxiv.org/abs/2609.07169v1. Full licence notice: \texttt{licenses/arxiv-2609.07169-CC-BY-4.0.txt}.\par\smallskip

\noindent\textit{Editorial scope.} Counted unit: the article's theorem thm:lipschitz-payoff-convergence-rate (source0.tex lines 469--486). The article's complete original proof of it is delivered with it (source0.tex lines 1787--1855, one \texttt{proof} environment of 69 lines, which the article places away from the statement). No mathematics is written, repaired or re-derived: the statement and the proof are the article's own lines, in the article's own notation, and no retained line is substituted or re-flowed.  Retained uncounted as the source-local context the proof needs: lines 437--439.  Nothing was omitted from any retained span, so the package carries no omission record.  No retained line contains a citation, so the wrapper carries no bibliography and no bibliography entry.  No retained line contains a figure environment, an image-inclusion command or drawing code, so no image is delivered and the package declares no asset dependency.  Editorial formatting only, in addition: an AMS \texttt{amsart} preamble that restates the article's own declarations and macros used by the retained lines, namely the environments \texttt{theorem} on separate counters, together with the standard packages those lines need.  The retained environments print in this document as Theorem 1 in order of appearance, whereas the article prints the same items with the numbering of its own full document.  The complete native source of the article as distributed in the arXiv e-print for this exact version is bundled unchanged as \texttt{sources/209606/source0.tex}.
\begin{align}\label{eq:gN-delta-general}
    g_N^\Delta(\mathrm{d}t) := \frac{1}{T_N}\sum_{j=0}^N\omega_j^N\delta_{t_j}(\mathrm{d}t).
\end{align}

\begin{theorem}\label{thm:lipschitz-payoff-convergence-rate}
    Let $(\mathcal P_T^N)_{N\geq1}$ be a sequence of monitoring grids with mesh size $\Delta_N\to0$, and let $g_N^\Delta$ be defined by~\eqref{eq:gN-delta-general}. Suppose that $\exists q>1$ such that
    \[
        \quad \mathbb E[S_T^{2q}]<\infty.
    \]
    Set $X_t:=\log S_t$ and suppose that $X$ admits the semimartingale decomposition $X_t=X_0+A_t+M_t$, where $A$ is a continuous finite-variation process with $A_0=0$ and $M$ is a continuous local martingale with $M_0=0$, such that
    \begin{align}\label{eq:log-semimartingale-rate-assumption}
        \big\||A|_T\big\|_{L^{2q}(\Omega)} + \big\|\langle M\rangle_T^{1/2}\big\|_{L^{2q}(\Omega)} < \infty.
    \end{align}
    Let $h:\mathbb R_+^2\longrightarrow \mathbb R$ satisfy
    \begin{align}\label{eq:lipschitz-first-variable}
        |h(x,y)-h(x',y)| \leq L_h|x-x'|, \qquad x,x',y>0,
    \end{align}
    for some $L_h \in (0,\infty)$, and assume that $h(1,S_T)\in L^q(\Omega)$. Then there exists a constant $\Lambda_{q,h}(T)>0$, independent of $N$, such that
    \begin{align}\label{eq:lipschitz-price-process-rate}
        \mathbb E\left[ \sup_{t\in[0,T]} |H_t^{h,N}-H_t^{h,c}|^q \right] \leq \Lambda_{q,h}(T)\Delta_N^q.
    \end{align}
\end{theorem}

\begin{proof}[Proof of Theorem \ref{thm:lipschitz-payoff-convergence-rate}]
    Set $S_T^* := \sup_{u\in[0,T]}S_u$. Since $(\mathrm e^{-rt}S_t)_{t\in[0,T]}$ is a martingale, $r\geq0$, and $S_T\in L^{2q}(\Omega)$, Doob's inequality gives $\|S_T^*\|_{L^{2q}(\Omega)}<\infty$. Note also that the prices are well defined. Indeed, by
    \eqref{eq:lipschitz-first-variable},
    \[
        |h(G_T^N,S_T)| \leq |h(1,S_T)| + L_h|G_T^N-1| \leq |h(1,S_T)| + L_h(1+S_T^*),
    \]
    which belongs to $L^q(\Omega)$. The same argument applies to $h(G_T^c,S_T)$.
    
    Set $L_N:=\log G_T^N$, $L_c:=\log G_T^c$, $\delta_N:=\omega_N^N=t_N-t_{N-1}$ and $T_N=T+\delta_N$. Then $\delta_N\leq\Delta_N$. Define the left-Riemann approximation
    \[
        \bar L_N := \frac1T \sum_{j=0}^{N-1} \omega_j^N X_{t_j}.
    \]
    By the definition of $G_T^N$, $L_N = \frac{T}{T_N}\bar L_N + \frac{\delta_N}{T_N}X_T$. We first estimate $\bar L_N-L_c$. Since $L_c = \frac1T\int_0^T X_s\,\mathrm ds$, we have 
    \begin{align*}
        T(\bar L_N-L_c) &= \sum_{j=0}^{N-1} \int_{t_j}^{t_{j+1}} (X_{t_j}-X_s)\,\mathrm ds.
    \end{align*}
    For $u\in(t_j,t_{j+1}]$, define $\kappa_N(u):=t_{j+1}-u$, such that $0\leq\kappa_N(u)\leq\Delta_N$. Integration by parts on each interval gives
    \begin{align}\label{eq:log-riemann-error}
        T(\bar L_N-L_c) &= -\sum_{j=0}^{N-1} \int_{t_j}^{t_{j+1}} (t_{j+1}-u)\,\mathrm dX_u
        = -\int_0^T\kappa_N(u)\,\mathrm dA_u - \int_0^T\kappa_N(u)\,\mathrm dM_u.
    \end{align}
    Since $A$ has finite variation, 
    \begin{align}\label{eq:A-riemann-bound}
        \left\| \int_0^T\kappa_N(u)\,\mathrm dA_u \right\|_{L^{2q}(\Omega)} \leq \Delta_N \big\||A|_T\big\|_{L^{2q}(\Omega)}.
    \end{align}
    By the Burkholder--Davis--Gundy inequality,
    \begin{align*}
        \left\| \int_0^T\kappa_N(u)\,\mathrm dM_u \right\|_{L^{2q}(\Omega)} \leq C_{2q} \left\| \left( \int_0^T \kappa_N(u)^2\, \mathrm d\langle M\rangle_u \right)^{1/2} \right\|_{L^{2q}(\Omega)} \leq C_{2q}\Delta_N \big\| \langle M\rangle_T^{1/2} \big\|_{L^{2q}(\Omega)}.
    \end{align*}
    Therefore
    \begin{align}\label{eq:left-riemann-log-rate}
        \|\bar L_N-L_c\|_{L^{2q}(\Omega)} \leq \frac{\Delta_N}{T} \left( \big\||A|_T\big\|_{L^{2q}(\Omega)} + C_{2q} \big\|\langle M\rangle_T^{1/2}\big\|_{L^{2q}(\Omega)} \right).
    \end{align}

    It remains to control the remaining terminal observation point. Let us write 
    \begin{align}\label{eq:LN-minus-Lc}
        L_N-L_c &= \frac{T}{T_N}(\bar L_N-L_c) + \frac{\delta_N}{T_N}(X_T-L_c).
    \end{align}
    By integration by parts,
    \begin{align*}
        X_T-L_c = \frac1T \left( TX_T-\int_0^TX_s\,\mathrm ds \right)
        &= \frac1T\int_0^T u\,\mathrm dX_u 
        \\ &= \frac1T\int_0^T u\,\mathrm dA_u + \frac1T\int_0^T u\,\mathrm dM_u.
    \end{align*}
    Using the BDG-inequality we arrive at $\|X_T-L_c\|_{L^{2q}(\Omega)} \leq \big\||A|_T\big\|_{L^{2q}(\Omega)} + C_{2q} \big\|\langle M\rangle_T^{1/2}\big\|_{L^{2q}(\Omega)}$. This, combined with \eqref{eq:left-riemann-log-rate}, \eqref{eq:LN-minus-Lc}, and using $\frac{T}{T_N}\leq1$ and $\frac{\delta_N}{T_N} \leq \frac{\Delta_N}{T}$, yields
    \begin{align}\label{eq:log-average-rate}
        \|L_N-L_c\|_{L^{2q}(\Omega)} \leq \frac{2\Delta_N}{T} \left( \big\||A|_T\big\|_{L^{2q}(\Omega)} + C_{2q} \big\|\langle M\rangle_T^{1/2}\big\|_{L^{2q}(\Omega)} \right).
    \end{align}
    
    We next pass from the logarithmic averages to the geometric averages. By Jensen's inequality,
    \[
        G_T^N \leq \int_{[0,T]}S_u\,g_N^\Delta(\mathrm du) \leq S_T^*\quad \text{ and } \quad G_T^c \leq \frac1T\int_0^T S_u\,\mathrm du \leq S_T^*.
    \]
    Using $ |\mathrm e^x-\mathrm e^y| \leq (\mathrm e^x+\mathrm e^y)|x-y|$ for $x,y\in\mathbb R$, we obtain 
    \begin{align*}
        |G_T^N-G_T^c| \leq (G_T^N+G_T^c)|L_N-L_c| \leq 2S_T^*|L_N-L_c|.
    \end{align*}
    By the Lipschitz property~\eqref{eq:lipschitz-first-variable}, H\"older's inequality, and~\eqref{eq:log-average-rate}, we obtain
    \begin{align*}
        \|h(G_T^N,S_T)-h(G_T^c,S_T)\|_{L^q(\Omega)} &\leq L_h \|G_T^N-G_T^c\|_{L^q(\Omega)} 
        \\ &\leq 2L_h \|S_T^*\|_{L^{2q}(\Omega)}\|L_N-L_c\|_{L^{2q}(\Omega)} 
        \leq 4C_q(T) \Delta_N \|S_T^*\|_{L^{2q}(\Omega)}
    \end{align*}
    for a finite constant $C_q(T)$ independent of $N$. Finally, recalling that $D_N := h(G_T^N,S_T)-h(G_T^c,S_T)$. Then $H_t^{h,N}-H_t^{h,c} = \mathrm e^{-r(T-t)} \mathbb E[D_N\mid\mathcal F_t]$. Since $r\geq0$, $\mathrm e^{-r(T-t)}\leq1$, and therefore $\sup_{t\in[0,T]} |H_t^{h,N}-H_t^{h,c}| \leq \sup_{t\in[0,T]} \left| \mathbb E[D_N\mid\mathcal F_t] \right|$. By Doob's $L^q$ inequality,
    \begin{align*}
        \left\| \sup_{t\in[0,T]} |H_t^{h,N}-H_t^{h,c}| \right\|_{L^q(\Omega)} \leq \frac{q}{q-1} \|D_N\|_{L^q(\Omega)} \leq \frac{q}{q-1} L_h C_q(T)\Delta_N.
    \end{align*}
    Taking the $q$-th power yields \eqref{eq:lipschitz-price-process-rate}.
\end{proof}

\end{document}
